% =============================================================================
% Offline Adjoint Validation for GCHP 4D-Var: Technical Description
% =============================================================================
\documentclass[12pt]{article}

\usepackage{amsmath,amssymb,amsfonts}
\usepackage{geometry}
\usepackage{hyperref}
\usepackage{booktabs}
\usepackage{array}
\usepackage{xcolor}
\usepackage{bm}
\usepackage{physics}
\usepackage{listings}
\usepackage{enumitem}

\geometry{margin=1in}

\newcommand{\sigb}{\sigma_b}
\newcommand{\gfd}{g^{\mathrm{FD}}}
\newcommand{\gadj}{g^{\mathrm{adj}}}
\newcommand{\gfit}{g^{\mathrm{fit}}}
\newcommand{\ek}{\mathbf{e}_k}
\newcommand{\eps}{\varepsilon}

\lstset{basicstyle=\small\ttfamily, breaklines=true,
        frame=single, backgroundcolor=\color{gray!8}}

\title{\textbf{Offline Adjoint Validation in GCHP 4D-Var}\\
\large Finite-Difference Gradient Test on the Native Cubed-Sphere Grid}
\author{GCHP Adjoint Development}
\date{\today}

% =============================================================================
\begin{document}
\maketitle
\tableofcontents
\newpage

% =============================================================================
\section{Overview}
% =============================================================================

The adjoint model must satisfy the identity
\begin{equation}
  \langle \mathbf{M}\,\mathbf{u},\,\mathbf{v} \rangle
  = \langle \mathbf{u},\,\mathbf{M}^{\top}\mathbf{v} \rangle
\end{equation}
for all vectors $\mathbf{u}$, $\mathbf{v}$, where $\mathbf{M}$ is the
tangent-linear model and $\mathbf{M}^{\top}$ is its adjoint.  In practice,
the adjoint is coded by hand and errors in the adjoint chain (incorrect
accumulation, missing terms, sign errors) produce a gradient
$\gadj = \partial J/\partial \sigma$ that disagrees with the true gradient.

This document describes the \emph{offline} validation framework implemented in
\texttt{test\_adjoint\_vs\_forward.run} and \texttt{cs\_fd\_tools.py}.  The
approach constructs the exact finite-difference (FD) gradient from perturbed
forward runs, then compares it cell by cell with the adjoint gradient.  Because
the cost function is \emph{exactly quadratic} in the control variable $\sigma$
under the conditions of the test (linear transport with no flux limiter), the
symmetric FD estimate is exact --- any discrepancy is a true adjoint bug, not a
FD approximation error.

All perturbations and gradient comparisons are performed on the \emph{native
cubed-sphere} (CS) grid.  This eliminates both the lat-lon-to-CS emission
regrid and the CS-to-lat-lon gradient regrid from the comparison, isolating
errors in the adjoint chain itself.

% =============================================================================
\section{Mathematical Framework}
% =============================================================================

\subsection{Forward Model}

Let $\sigma(f,j,i) \in \mathbb{R}$ be the dimensionless emission scaling
factor at cubed-sphere cell $(f,j,i)$, where $f\in\{0,\ldots,5\}$ is the face
index and $(j,i)$ are the within-face row and column.  The prior TER flux
$F_{\mathrm{base}}(f,j,i)$ is fixed; the actual flux applied by HEMCO is
\begin{equation}
  F_{\mathrm{total}}(f,j,i) = \sigma(f,j,i)\cdot F_{\mathrm{base}}(f,j,i).
\end{equation}

The forward model $\mathcal{M}$ integrates the CO$_2$ transport equations
forced by $F_{\mathrm{total}}$ over the assimilation window
$[t_0, T]$:
\begin{equation}
  \mathbf{x}(t) = \mathcal{M}_{t_0 \to t}(\mathbf{x}_0,\,\sigma).
\end{equation}
With chemistry, convection, and PBL mixing disabled and the PPM advection
limiter switched off (\texttt{hord\_tr=2}), $\mathcal{M}$ is \emph{exactly
linear} in $\sigma$:
\begin{equation}
  \mathbf{x}(t;\sigma) = \mathbf{x}^{(0)}(t) + \mathbf{A}(t)\,\bm{\sigma},
  \label{eq:linearity}
\end{equation}
where $\mathbf{x}^{(0)}$ is the state driven by all non-TER fluxes,
$\bm{\sigma}$ is the vectorised scaling field, and $\mathbf{A}$ is the linear
transport operator mapping emission scalings to CO$_2$ concentrations.

\subsubsection*{Role of the PPM limiter}

The default advection scheme (\texttt{hord\_tr=12}) applies a monotone PPM
limiter that is a non-linear function of the tracer concentration.  This
renders $\mathcal{M}$ non-linear in $\sigma$, breaking the quadraticity
argument below.  By setting \texttt{hord\_tr=2} (unlimited PPM), the transport
operator becomes exactly linear in the tracer and hence in $\sigma$, restoring
quadraticity.

\subsection{Cost Function}

The observation term of the cost function is
\begin{equation}
  J(\sigma) = \frac{1}{2} \sum_{k=1}^{N_{\mathrm{obs}}}
  \left(\frac{\hat{x}_{\mathrm{XCO}_2,k}(\sigma) - y_k}{\epsilon_k}
  \right)^2,
  \label{eq:J}
\end{equation}
where $\hat{x}_{\mathrm{XCO}_2,k}(\sigma)$ is the simulated column-averaged
CO$_2$ at observation $k$, $y_k$ is the synthetic observation, and $\epsilon_k$
is the observation error standard deviation.

By \eqref{eq:linearity}, $\hat{x}_{\mathrm{XCO}_2,k}$ is an affine function of
$\bm{\sigma}$:
\begin{equation}
  \hat{x}_{\mathrm{XCO}_2,k}(\sigma) = c_k + \mathbf{b}_k^\top\bm{\sigma},
\end{equation}
so $J(\sigma)$ is a sum of squares of affine functions of $\bm{\sigma}$ and is
therefore \emph{exactly quadratic} in $\bm{\sigma}$:
\begin{equation}
  J(\bm{\sigma}) = \frac{1}{2}\bm{\sigma}^\top \mathbf{Q}\,\bm{\sigma}
                   + \mathbf{p}^\top\bm{\sigma} + c_0,
  \label{eq:J_quadratic}
\end{equation}
where $\mathbf{Q} = \sum_k \mathbf{b}_k\mathbf{b}_k^\top / \epsilon_k^2$
is the Gauss-Newton Hessian.

\subsection{Adjoint Gradient}

The adjoint model computes the exact gradient
\begin{equation}
  \gadj(f,j,i) = \frac{\partial J}{\partial \sigma(f,j,i)}.
  \label{eq:g_adj}
\end{equation}
In the GCHP adjoint, this is accumulated into the field
\texttt{SurfaceFluxAdj\_CO2} written to the adjoint \texttt{OutputDir}.

\subsection{Finite-Difference Gradient and Exactness}

For a unit perturbation along cell $k = (f,j,i)$ define
\begin{align}
  \sigma^{+} &= \sigma_0 + \eps\,\ek, \\
  \sigma^{-} &= \sigma_0 - \eps\,\ek,
\end{align}
where $\ek$ is the standard basis vector for cell $k$ and $\eps > 0$ is the
perturbation amplitude.  The symmetric (centred) finite-difference estimate is
\begin{equation}
  \gfd_k(\eps) = \frac{J(\sigma^{+}) - J(\sigma^{-})}{2\eps}.
  \label{eq:gFD}
\end{equation}

Because $J$ is exactly quadratic \eqref{eq:J_quadratic}, its Taylor expansion
terminates at second order:
\begin{align}
  J(\sigma^{+}) &= J(\sigma_0) + \eps\,\gadj_k
                   + \tfrac{1}{2}\eps^2\,Q_{kk}, \\
  J(\sigma^{-}) &= J(\sigma_0) - \eps\,\gadj_k
                   + \tfrac{1}{2}\eps^2\,Q_{kk}.
\end{align}
Substituting into \eqref{eq:gFD}:
\begin{equation}
  \boxed{\gfd_k(\eps) = \gadj_k \quad \text{exactly, for any } \eps > 0.}
  \label{eq:exactness}
\end{equation}
There is \emph{no} finite-difference truncation error.  Any discrepancy between
$\gfd_k$ and $\gadj_k$ is therefore a genuine error in the adjoint code.

\subsection{Parabola Fit and the $g^{\mathrm{fit}}$ Estimator}

In practice, rounding errors in floating-point arithmetic introduce small noise.
To maximally exploit all available $J$ evaluations, a parabola
\begin{equation}
  J(\eps) = J_0 + \eps\,\gfit_k + \tfrac{1}{2}\eps^2\,H_{kk}
  \label{eq:parabola}
\end{equation}
is fitted by least squares through the base point $(0, J_0)$ and all
perturbed evaluations $\{(\eps_n, J(\sigma_0 + \eps_n\,\ek))\}$.  The fit
yields $\gfit_k$ and the diagonal Hessian element $H_{kk}$; the residual of
the fit, $\max_n|\,J_{\mathrm{obs},n} - J_{\mathrm{fit},n}|$, provides a
diagnostic for the degree of quadraticity.

The validation metric is
\begin{equation}
  r_k = \frac{\gadj_k}{\gfit_k}.
\end{equation}
For a correct adjoint, $r_k = 1$ to within floating-point precision.

\subsection{Global Perturbation Test}

In addition to cell-wise perturbations, a \emph{global} uniform perturbation
$\sigma \to (1+\eps)\sigma_0$ is applied.  For this perturbation the
expected gradient from the adjoint is
\begin{equation}
  g_{\mathrm{global}}^{\mathrm{adj}} =
  \sum_{f,j,i} \sigma_0(f,j,i)\cdot\gadj(f,j,i),
  \label{eq:g_global}
\end{equation}
which is compared against
$\gfd_{\mathrm{global}} = [J((1+\eps)\sigma_0) - J((1-\eps)\sigma_0)]/(2\eps)$.
This test is sensitive to errors that affect many cells simultaneously, such as
a missing overall scale factor.

\subsection{Quadraticity Probe}

For the top-ranked cell (largest $|\gadj_k|$), an additional one-sided
evaluation at $+2\eps$ is collected:
\begin{equation}
  J(\sigma_0 + 2\eps\,\ek) = J_0 + 2\eps\,\gadj_k + 2\eps^2\,H_{kk}.
\end{equation}
Together with $J(\sigma_0 \pm \eps\,\ek)$, this gives three points to fit the
parabola \eqref{eq:parabola}.  If the residual of the three-point fit is small
relative to $\max_n |J_n - J_0|$ (threshold $10^{-3}$), $J$ is confirmed to
be quadratic; otherwise a non-linearity is flagged, which would indicate that
the PPM limiter or chemistry is inadvertently active.

% =============================================================================
\section{Native Cubed-Sphere Validation Strategy}
% =============================================================================

\subsection{Motivation}

An earlier lat-lon FD test (fd\_multicell) revealed a gradient mismatch that
could stem from two sources:
\begin{enumerate}[label=(\roman*)]
  \item An error in the adjoint chain.
  \item A mismatch between the forward emission regrid (lat-lon $\to$ CS,
        operator $\mathbf{S}$) and the adjoint gradient regrid (CS $\to$
        lat-lon, operator $\mathbf{R}$), i.e.\ $\mathbf{S} \neq \mathbf{R}^\top$.
\end{enumerate}
By placing $\sigma$ directly on the native CS grid and outputting the adjoint
gradient on the same native CS grid (no \texttt{grid\_label} in \texttt{HISTORY.rc}),
both regrid operators are removed from the loop.  A mismatch on the native
grid is unambiguously a bug in the adjoint chain, not a regridding inconsistency.

\subsection{CS Sigma File Format}

The sigma field is written as a NetCDF4 file with dimensions
\texttt{(time, lev=1, lat=$6N$, lon=$N$)}, where $N$ is the CS resolution
(\texttt{IM}, default 24 for C24).  The $6N$ rows correspond to the six faces
stacked vertically; MAPL ExtData detects the aspect ratio and reads the file
as a cubed-sphere field with an identity regrid (no interpolation).

\subsection{P0 Plumbing Gate}

Before any expensive month-long run, a one-day forward pair verifies that
ExtData delivers the CS sigma to HEMCO without regridding:
\begin{enumerate}
  \item \textbf{Base run}: $\sigma \equiv 1$ everywhere; record
        \texttt{EmisCO2\_NetTerrExch} output $E^{(0)}_{f,j,i}$.
  \item \textbf{Unique-per-cell run}: assign a unique value
        $\sigma^{\mathrm{uniq}}_{f,j,i}$ to every cell; record
        $E^{(1)}_{f,j,i}$.
  \item \textbf{Check}: verify $E^{(1)}_{f,j,i} = \sigma^{\mathrm{uniq}}_{f,j,i}
        \cdot E^{(0)}_{f,j,i}$ to machine precision at every cell.
\end{enumerate}
If the check fails, ExtData has regridded the CS file (breaking the
native-grid assumption) or the face mapping is wrong; the script aborts
immediately.

% =============================================================================
\section{Implementation}
% =============================================================================

The test is orchestrated by the PBS script
\texttt{test\_adjoint\_vs\_forward.run} (submitted from the test directory)
and the Python utility \texttt{cs\_fd\_tools.py}.  The run is divided into
five phases; marker files in \texttt{results/} allow resumption
(\texttt{RESTART=true}) if the job is interrupted.

\subsection{Phase P0: CS Plumbing Gate (two 1-day forwards)}

\begin{enumerate}
  \item \textbf{Write sigma} with \texttt{--kind none} (uniform 1.0).
  \item Run GCHP forward for one day (\texttt{T\_START} to
        \texttt{T\_START+1d}).
  \item \textbf{Write sigma} with \texttt{--kind unique} (unique per-cell
        pattern of amplitude 0.5 around 1.0).
  \item Run GCHP forward for one day.
  \item \textbf{Check} (Python \texttt{check-p0} subcommand): assert
        $E^{(1)} = \sigma^{\mathrm{uniq}} \cdot E^{(0)}$ at every CS cell.
        If the check fails with \texttt{die}, the run aborts.
\end{enumerate}

\subsection{Phase P1: Base Forward Run and Adjoint Forcing}

\begin{enumerate}
  \item \textbf{Write sigma} with \texttt{--kind none} ($\sigma \equiv 1$).
  \item Run GCHP forward over the full assimilation window
        $[t_0, T]$, producing:
        \begin{itemize}
          \item Adjoint checkpoints (\texttt{gcadj\_internal\_checkpoint*},
                \texttt{gcadj\_import\_checkpoint*}) in \texttt{forward\_run/}.
          \item Chemistry restart
                \texttt{Restarts/gcchem\_internal\_checkpoint*}.
          \item Satellite-track diagnostic \texttt{GEOSChem.sat\_track.*.nc4}
                in \texttt{forward\_run/OutputDir/}.
        \end{itemize}
  \item Move checkpoints and restart to \texttt{adjoint/}:
        \begin{lstlisting}
find forward_run/ -name 'gcadj_*checkpoint*' \
     -exec mv -t adjoint/ {} +
mv forward_run/Restarts/gcchem_internal_checkpoint* \
     adjoint/gcchem_initial_restart.nc4
        \end{lstlisting}
  \item Compute $J_0$ and adjoint forcing from the sat-track files
        (\texttt{co2\_adjoint\_forcing\_osse.py}); write
        \texttt{CO2\_adjoint\_forcing\_*.nc4} to \texttt{forcing\_files\_osse/}
        and record $J_0$ in \texttt{results/J\_base.txt}.
\end{enumerate}

\subsection{Phase P2: Adjoint Run and Cell Selection}

\begin{enumerate}
  \item Run GCHP adjoint over $[t_0, T]$, reading checkpoints from
        \texttt{adjoint/} and forcing from \texttt{forcing\_files\_osse/}.
        The adjoint outputs \texttt{GEOSChem.SurfaceFluxAdj\_CO2.*.nc4}
        in \texttt{adjoint/OutputDir/} on the native C24 CS grid (no
        \texttt{grid\_label} in \texttt{adjoint/HISTORY.rc}).
  \item \textbf{Pick cells} (\texttt{cs\_fd\_tools.py pick-cells}):
        \begin{itemize}
          \item Accumulate the adjoint gradient field
                $\gadj(f,j,i) = \partial J / \partial \sigma(f,j,i)$
                over the window.
          \item Select the $N$ land cells
                (determined from a coordinate file) with the largest
                $|\gadj(f,j,i)| > g_{\min}$, ranked in descending order.
          \item Write \texttt{results/cells.txt}: one row per perturbation run,
                recording \texttt{(kind, f, j, i, eps, tag)}.
          \item For the top-ranked cell add an extra row at $+2\eps$
                (quadraticity probe).
          \item Add two global-perturbation rows ($\pm\eps$).
          \item Save the full gradient array to
                \texttt{results/cs\_state.npz}.
        \end{itemize}
  \item Delete adjoint checkpoints to free disk space.
\end{enumerate}

\subsection{Phase P3: Perturbed Forward Runs}

For each row in \texttt{results/cells.txt} (processed sequentially):
\begin{enumerate}
  \item \textbf{Write sigma}: for a cell perturbation, set
        $\sigma(f,j,i) = 1 + \eps$ at a single CS cell; for a global
        perturbation, $\sigma \equiv 1 + \eps$ everywhere.
  \item Run GCHP forward over $[t_0, T]$.
  \item Compute $J(\sigma^{+})$ or $J(\sigma^{-})$ with
        \texttt{co2\_adjoint\_forcing\_osse.py}; write
        \texttt{results/J\_\{tag\}.txt}.
  \item Checkpoints are dropped immediately (not needed for subsequent
        phases).
\end{enumerate}
A walltime guard checks remaining PBS time before each P3 run; if budget
is exhausted the run exits with code 5 and can be resubmitted with
\texttt{RESTART=true}.

\subsection{Phase P4: Analysis and Verdict}

\texttt{cs\_fd\_tools.py analyze} performs the following:

\begin{enumerate}
  \item For each cell group $(f,j,i)$, collect all $J$ evaluations
        $\{(\eps_n, J_n)\}$ and fit the parabola \eqref{eq:parabola} by
        least squares to obtain $\gfit_k$ and $H_{kk}$.
  \item Compute the ratio $r_k = \gadj_k / \gfit_k$ for each cell.
  \item Compute the global-perturbation ratio
        $r_{\mathrm{global}} = \gadj_{\mathrm{global}} / \gfit_{\mathrm{global}}$.
  \item Assess quadraticity: the parabola residual must satisfy
        $\max_n|\mathrm{resid}_n| < 10^{-3} \max_n|J_n - J_0|$.
  \item Write \texttt{results/report.txt}, scatter plot
        \texttt{results/fd\_scatter\_cs.png} (adjoint gradient vs.\ FD
        gradient), and spatial map \texttt{results/g\_map\_cs.png}.
  \item Return an exit code:
\end{enumerate}

\begin{table}[h!]
\centering
\caption{Exit codes and their interpretation.}
\begin{tabular}{cll}
\toprule
Code & Meaning & Action \\
\midrule
0 & Adjoint validated: $|r_k - 1| \leq \delta_{\mathrm{tol}}$ for all cells
  & None \\
3 & Systematic factor: $r_k \approx \alpha \neq 1$ (constant offset)
  & Check overall scale (unit conversion, area factor) \\
4 & Spatially varying ratio: $r_k$ varies across cells
  & Adjoint chain bug; inspect cell-by-cell scatter \\
5 & Incomplete: not all P3 runs finished in walltime
  & Resubmit with \texttt{RESTART=true} \\
\bottomrule
\end{tabular}
\end{table}

% =============================================================================
\section{Interpreting Results}
% =============================================================================

\subsection{Scatter Plot (\texttt{fd\_scatter\_cs.png})}

Each point plots $\gadj_k$ (y-axis) against $\gfit_k$ (x-axis).  For a
correct adjoint all points lie on the diagonal $y = x$.  A systematic offset
(all points on a line $y = \alpha x$, $\alpha \neq 1$) indicates a global
scale error.  Scatter around any line indicates spatially varying errors in the
adjoint chain.

\subsection{Spatial Map (\texttt{g\_map\_cs.png})}

Shows the full adjoint gradient field $\gadj(f,j,i)$ on the globe, with the
FD test cells marked and labelled with their ratio $r_k = \gadj_k/\gfit_k$.

\subsection{One-Sided FD Diagnostics}

The report also lists one-sided estimates
$[J(\sigma_0 + \eps\,\ek) - J_0] / \eps$, which include the $O(\eps)$ Hessian
contribution $\eps H_{kk}/2$.  Agreement between one-sided estimates at
different $\eps$ and the parabola fit confirms quadraticity.

% =============================================================================
\section{Configuration}
% =============================================================================

\subsection{Required Physics Settings}

The following settings in \texttt{runConfig\_forward.sh} and
\texttt{runConfig\_adjoint.sh} are required for the quadraticity argument to
hold:

\begin{table}[h!]
\centering
\begin{tabular}{ll}
\toprule
Setting & Required value \\
\midrule
\texttt{Turn\_on\_Chemistry}         & \texttt{false} \\
\texttt{Turn\_on\_Cloud\_Conv}       & \texttt{false} \\
\texttt{Turn\_on\_PBL\_Mixing}       & \texttt{false} \\
\texttt{Turn\_on\_Non\_Local\_Mixing}& \texttt{false} \\
\texttt{Turn\_on\_Dry\_Deposition}   & \texttt{false} \\
\texttt{Turn\_on\_Wet\_Deposition}   & \texttt{false} \\
\texttt{Turn\_on\_Transport}         & \texttt{true} \\
\texttt{Turn\_on\_Base\_Emissions}   & \texttt{true} \\
\bottomrule
\end{tabular}
\end{table}

\textbf{Rationale.} Chemistry is a non-linear sink; convection and PBL mixing
introduce non-linear mixing operators; dry and wet deposition introduce
concentration-dependent sinks.  All of these break the linearity
\eqref{eq:linearity} of $\mathcal{M}$ in $\sigma$ and hence the quadraticity
of $J$.  With only transport and base emissions active, $\sigma$ enters $J$
solely through the linear transport of the surface emission pulse, preserving
exact quadraticity.

\subsection{Advection Limiter}

The \texttt{input.nml} must set

\begin{lstlisting}
&fv_core_nml
hord_tr = 2 /
\end{lstlisting}

\texttt{hord\_tr=2} selects unlimited PPM advection, which is exactly linear in
the tracer concentration.  The default \texttt{hord\_tr=12} applies a monotone
limiter that is non-linear in the tracer, breaking quadraticity.

\subsection{ExtData Offsets for the Adjoint}

The GCHP adjoint runs time-reversed: at each backward step the model reads
meteorology at time $t - \Delta t$ rather than $t$.  The MAPL ExtData refresh
offsets in \texttt{adjoint/ExtData.rc} must therefore be shifted by one
timestep $\Delta t$ relative to the forward values:

\begin{table}[h!]
\centering
\begin{tabular}{llll}
\toprule
Met group & Forward offset & Adjoint offset ($\Delta t = 10$\,min)
          & Adjoint offset ($\Delta t = 30$\,min) \\
\midrule
A1 (1-hr averaged 2-D) & \texttt{F0;003000} & \texttt{F0;002000} & \texttt{F0;000000} \\
A3 (3-hr averaged 3-D) & \texttt{F0;013000} & \texttt{F0;012000} & \texttt{F0;010000} \\
I3 first slot (PS1, SPHU1, TMPU1) & \texttt{0} & \texttt{0;-001000} & \texttt{0;-003000} \\
I3 second slot (PS2, SPHU2, TMPU2) & \texttt{0;003000} & \texttt{0} & \texttt{0} \\
\bottomrule
\end{tabular}
\end{table}

The general rule is: adjoint offset $=$ forward offset $-\,\Delta t$.

% =============================================================================
\section{Key Files}
% =============================================================================

\begin{table}[h!]
\centering
\begin{tabular}{ll}
\toprule
File & Purpose \\
\midrule
\texttt{test\_adjoint\_vs\_forward.run}
  & PBS script orchestrating all phases \\
\texttt{cs\_fd\_tools.py}
  & Python: write-sigma, check-p0, pick-cells, analyze \\
\texttt{co2\_adjoint\_forcing\_osse.py}
  & Compute $J$ and adjoint forcing from sat-track files \\
\texttt{forward\_run/ExtData.rc}
  & ExtData config for forward run (forward offsets) \\
\texttt{adjoint/ExtData.rc}
  & ExtData config for adjoint run (shifted offsets) \\
\texttt{adjoint/HISTORY.rc}
  & Must omit \texttt{grid\_label} for native-CS gradient output \\
\texttt{sigma\_monthly/\{year\}/\{mm\}.nc4}
  & CS sigma field written by \texttt{write-sigma} before each run \\
\texttt{results/cells.txt}
  & FD run list: one row per perturbed forward \\
\texttt{results/cs\_state.npz}
  & Adjoint gradient array and cell coordinates \\
\texttt{results/report.txt}
  & Per-cell FD vs.\ adjoint comparison and verdict \\
\texttt{results/fd\_scatter\_cs.png}
  & Scatter: $\gadj$ vs.\ $\gfit$ \\
\texttt{results/g\_map\_cs.png}
  & Spatial map of $\gadj$ with FD cells labelled \\
\bottomrule
\end{tabular}
\end{table}

% =============================================================================
\section{Relationship to the Full 4D-Var System}
% =============================================================================

This offline test validates the adjoint \emph{transport} chain in isolation.
Once validated, the full 4D-Var system adds:
\begin{itemize}
  \item The background term $J_b$ (computed analytically in Python, no GCHP
        involvement).
  \item The lat-lon $\to$ CS emission regrid (operator $\mathbf{S}$) and the
        CS $\to$ lat-lon gradient regrid (operator $\mathbf{R}$).
  \item Chemistry, convection, PBL mixing, and deposition (all of which are
        non-linear and require separate adjoint verification).
\end{itemize}
Passing this test with exit code 0 on the native CS grid guarantees that the
core transport adjoint is correct.  A subsequent lat-lon FD test (with the
same transport-only configuration but $\sigma$ on the lat-lon grid) would
then isolate any regrid inconsistency ($\mathbf{S} \neq \mathbf{R}^\top$).

\end{document}
